\section{Model}

Let output be produced by a constant-returns technology, so that
\begin{equation}
  Y_{t} = A_{t} K_{t}^{\alpha} L_{t}^{1 - \alpha},
  \label{eq:production}
\end{equation}
where $A_{t}$ is total factor productivity, $K_{t}$ is capital, and $L_{t}$ is
labour. The elasticity $\alpha$ is the parameter a reader is most likely to
argue with, which makes this a good paragraph to click on.

Equation~\eqref{eq:production} is deliberately unremarkable. What matters is
that this sentence lives in \texttt{sections/model.tex} rather than in
\texttt{main.tex}, so a click here exercises the path mapping back to the
project root.

\subsection{Households}

A representative household chooses consumption $C_{t}$ and next-period capital
$K_{t+1}$ to maximise
\begin{equation}
  \sum_{t = 0}^{\infty} \beta^{t} \, u(C_{t}),
  \qquad
  u(C) = \frac{C^{1 - \sigma} - 1}{1 - \sigma},
  \label{eq:utility}
\end{equation}
subject to the resource constraint
\begin{equation}
  C_{t} + K_{t+1} = Y_{t} + (1 - \delta) K_{t},
  \label{eq:budget}
\end{equation}
with discount factor $\beta \in (0, 1)$, intertemporal elasticity $1 / \sigma$,
and depreciation rate $\delta$. Labour is supplied inelastically at $L_{t} = 1$,
which is the sort of assumption that survives in a draft precisely because
nobody bothers to click on it.

Combining the first-order conditions for~\eqref{eq:utility} and~\eqref{eq:budget}
means using the budget constraint to write
$C_{t} = Y_{t} + (1 - \delta) K_{t} - K_{t+1}$, substituting it into the
objective, and differentiating with respect to the single remaining choice
variable $K_{t+1}$. Saving one more unit of capital costs one unit of consumption
today, which is worth $\beta^{t} u'(C_{t})$. Tomorrow that unit pays back two
things: the extra output it produces, $\alpha A_{t+1} K_{t+1}^{\alpha - 1}$,
which is the marginal product implied by~\eqref{eq:production}, and the
$1 - \delta$ units of capital left over after depreciation. Its gross return is
therefore $\alpha A_{t+1} K_{t+1}^{\alpha - 1} + 1 - \delta$, and the payoff is
worth $\beta^{t+1} u'(C_{t+1})$ times that return. Setting the cost equal to the
payoff and using $u'(C) = C^{-\sigma}$ gives the Euler equation
\begin{equation}
  \left( \frac{C_{t+1}}{C_{t}} \right)^{\sigma}
    = \beta \left[ \alpha A_{t+1} K_{t+1}^{\alpha - 1} + 1 - \delta \right],
  \label{eq:euler}
\end{equation}
which says that the utility the household gives up by saving one more unit
today must exactly match the discounted utility that unit returns tomorrow.
Everything that follows is a consequence of~\eqref{eq:euler} and a
transversality condition we will not write down.

\subsection{Steady state}

Setting $A_{t} = A$ and $C_{t+1} = C_{t}$ in~\eqref{eq:euler} pins down the
capital stock,
\begin{equation}
  K^{\star}
    = \left( \frac{\alpha A}{\beta^{-1} - 1 + \delta} \right)^{\frac{1}{1 - \alpha}},
  \label{eq:kstar}
\end{equation}
and steady-state consumption follows from~\eqref{eq:budget} as
$C^{\star} = A (K^{\star})^{\alpha} - \delta K^{\star}$. The comparative statics
are the ones every reader expects: $K^{\star}$ rises with patience and with
productivity, and falls as capital wears out faster.

Two features of~\eqref{eq:kstar} matter later. First, the exponent
$1 / (1 - \alpha)$ means that a small change in $\alpha$ moves the steady state a
lot, so the calibration of $\alpha$ deserves more attention than it usually
receives. Second, nothing in the steady state depends on $\sigma$; the
intertemporal elasticity governs the speed of transition rather than the
destination. Readers who dislike this separation tend to say so in the margin,
which is exactly the kind of comment this document is built to collect.
